\documentclass[letterpaper,11pt]{moderncv}
\moderncvstyle{banking}
\moderncvcolor{black}

\usepackage[letterpaper, margin=0.8in]{geometry}
\usepackage{xcolor}
\usepackage{fontspec}
\setmainfont{Times New Roman}
\usepackage[unicode]{hyperref}

\name{Saeyoung}{Kim}
\address{Cambridge, MA, USA}{}{}
\phone[mobile]{+1 (351) 322 5510}
\email{saeyoung\_kim@uml.edu}
\social[linkedin]{saeyoung-macx-kim}

\recipient{Hiring Team}{Anduril Industries\\Quincy, MA}
\date{\today}
\opening{Dear Hiring Team,}
\closing{Sincerely,}

\begin{document}

\makelettertitle

I'm applying for the Manufacturing Test Engineer role in Quincy. I've spent the last several years building test infrastructure for electromechanical hardware, and this role sits right at the intersection of what I've done and what I want to keep doing.

At Hanwha Systems, I developed and executed system-level test plans for military satellite programs. I designed EGSE test configurations, wrote controlled procedures that technicians could follow independently, performed root cause analysis when hardware failed during integration, and coordinated across electrical, mechanical, and software teams to close out issues. The pace was fast, the deadlines were real, and everything had to be documented and traceable.

During my PhD, I built the test side of things from scratch. I designed sensor hardware (circuit design, PCB layout, MEMS cleanroom fab), wrote firmware in C/C++ to talk to everything over UART, SPI, and I2C, then built test setups with oscilloscopes and function generators to characterize performance and catch failure modes. When something didn't work, I was the one figuring out whether it was a board-level issue, a firmware bug, or a process problem.

More recently, I've been writing Python tools for automated process monitoring and anomaly detection in a manufacturing context, which has sharpened my thinking around scalable, data-driven test approaches.

I'm based in Cambridge, so Quincy is an easy commute. I'd be glad to talk more about how I can help build out your test capabilities. Thanks for reading.

\makeletterclosing

\end{document}
